\documentclass[11pt,a4paper]{article}

\usepackage[utf8]{inputenc}
\usepackage[T1]{fontenc}
\usepackage[french]{babel}
\usepackage{lmodern}
\usepackage{geometry}
\usepackage{amsmath,amssymb,amsfonts,mathtools}
\usepackage{booktabs}
\usepackage{array}
\usepackage{longtable}
\usepackage{enumitem}
\usepackage{xcolor}
\usepackage{hyperref}
\usepackage{microtype}

\geometry{margin=2.4cm}
\hypersetup{
  colorlinks=true,
  linkcolor=blue!60!black,
  citecolor=blue!60!black,
  urlcolor=blue!60!black,
  pdftitle={Theorie holographique UV/IR a cellule de phase deformee},
  pdfauthor={Document de travail}
}

\setlength{\parindent}{0pt}
\setlength{\parskip}{0.55em}
\setlist[itemize]{topsep=3pt,itemsep=2pt,parsep=2pt}
\setlist[enumerate]{topsep=3pt,itemsep=2pt,parsep=2pt}

\newcommand{\lp}{\ell_{\rm P}}
\newcommand{\rp}{\rho_{\rm P}}
\newcommand{\rhoreg}{\rho_{\rm reg}}
\newcommand{\rhograv}{\rho_{\rm grav}}
\newcommand{\rholambda}{\rho_\Lambda}
\newcommand{\rhc}{\rho_c}
\newcommand{\mpl}{m_{\rm P}}
\newcommand{\Ep}{E_{\rm P}}
\newcommand{\Eh}{E_{\rm H}}
\newcommand{\RL}{R_\Lambda}
\newcommand{\LH}{L_H}
\newcommand{\calL}{\mathcal L}
\newcommand{\kb}{k_{\rm B}}
\newcommand{\dif}{\mathrm d}

\title{\textbf{Theorie holographique UV/IR a cellule de phase deformee}\\
\large Architecture, equations maitresses, no-go, projection holographique et lien possible avec $\alpha$}
\author{Document de travail}
\date{22 mai 2026}

\begin{document}
\maketitle

\begin{abstract}
Ce document rassemble une formulation coherente d'un cadre theorique effectif :
une cellule de phase deformee de type GUP regularise les integrales
ultraviolettes, mais cette regularisation locale ne suffit pas a expliquer
l'energie noire. Le resultat central est un no-go conditionnel : identifier
directement l'energie du vide regularisee a l'energie noire force un
\(\beta_0\sim10^{60}\), exclu comme parametre local universel. Une theorie
viable doit donc combiner la regularisation UV avec une projection
gravitationnelle IR de type holographique. Le meme rapport Planck--horizon
peut ensuite motiver une conjecture reliant \(\alpha\) a \(R_\Lambda\), dont le
verrou est le calcul du mode de bord Maxwell.
\end{abstract}

\tableofcontents
\newpage

\section{Statut epistemique}

Le cadre n'est pas une theorie fondamentale confirmee. Il separe quatre
niveaux :

\begin{center}
\begin{tabular}{@{}p{4cm}p{10cm}@{}}
\toprule
Niveau & Sens \\
\midrule
Resultat mathematique & consequence derivee dans la mesure GUP choisie. \\
Identite exacte & relation decoulant des definitions de \(c,\hbar,G,H_0,\Omega_\Lambda\). \\
No-go conditionnel & exclusion obtenue si l'on impose une interpretation directe. \\
Conjecture & relation structuree mais non encore derivee dynamiquement. \\
\bottomrule
\end{tabular}
\end{center}

La phrase directrice est :
\begin{equation}
\boxed{\text{le GUP regularise localement ; l'horizon projette globalement.}}
\end{equation}

\section{Definitions fondamentales}

\begin{equation}
\lp=\sqrt{\frac{\hbar G}{c^3}},
\qquad
\mpl=\sqrt{\frac{\hbar c}{G}},
\qquad
\Ep=\sqrt{\frac{\hbar c^5}{G}},
\end{equation}
\begin{equation}
\rp=\frac{c^7}{\hbar G^2},
\qquad
\LH=\frac{c}{H_0},
\qquad
\rhc=\frac{3c^2H_0^2}{8\pi G}.
\end{equation}

La densite d'energie noire observee est
\begin{equation}
\rho_{\rm DE}=\Omega_\Lambda \rhc.
\end{equation}

Pour un univers asymptotiquement de Sitter,
\begin{equation}
\RL=\sqrt{\frac{3}{\Lambda}}.
\end{equation}

L'entropie d'un horizon de rayon \(L\) est
\begin{equation}
S(L)=\frac{\kb A}{4\lp^2}=\frac{\kb\pi L^2}{\lp^2},
\qquad
N_{\rm holo}(L)=\frac{S(L)}{\kb}=\frac{\pi L^2}{\lp^2}.
\end{equation}

\section{Postulat I : cellule locale UV deformee}

On postule la mesure de phase
\begin{equation}
\boxed{
\dif N_\beta=
\frac{g\,\dif^3x\,\dif^3p}{h^3(1+\beta p^2)^\gamma}.
}
\end{equation}

Le cas central est
\begin{equation}
\boxed{
\gamma=3,
\qquad
\dif N_\beta=
\frac{g\,\dif^3x\,\dif^3p}{h^3(1+\beta p^2)^3}.
}
\end{equation}

Avec
\begin{equation}
\beta=\beta_0\frac{\lp^2}{\hbar^2}
=\frac{\beta_0}{\mpl^2c^2}.
\end{equation}

L'integrale maitresse est
\begin{equation}
I_s(\gamma)=\int_0^\infty \frac{p^{s-1}\,\dif p}{(1+\beta p^2)^\gamma}
=\frac12\beta^{-s/2}B\left(\frac{s}{2},\gamma-\frac{s}{2}\right),
\qquad \gamma>\frac{s}{2}.
\end{equation}

Pour \(\gamma=3\),
\begin{equation}
\int_0^\infty \frac{p^2\,\dif p}{(1+\beta p^2)^3}
=\frac{\pi}{16\beta^{3/2}},
\qquad
\int_0^\infty \frac{p^3\,\dif p}{(1+\beta p^2)^3}
=\frac{1}{4\beta^2}.
\end{equation}

\subsection{Densite maximale de modes}

\begin{equation}
\frac{N_{\max}}{V}
=\frac{4\pi g}{h^3}\int_0^\infty \frac{p^2\,\dif p}{(1+\beta p^2)^3}.
\end{equation}

Donc
\begin{equation}
\boxed{
n_{\max}=\frac{g\pi^2}{4h^3\beta^{3/2}}
=\frac{g}{32\pi\beta_0^{3/2}\lp^3}.
}
\end{equation}

\section{Postulat II : regularisation UV de l'energie du vide}

Pour un champ massless de degenerescence \(g\),
\begin{equation}
\rhoreg=\frac{4\pi g}{h^3}\int_0^\infty
\frac{p^2\,\dif p}{(1+\beta p^2)^3}\frac12 pc.
\end{equation}

Donc
\begin{equation}
\rhoreg
=\frac{g\pi c}{2h^3\beta^2}
=\boxed{\frac{g\rp}{16\pi^2\beta_0^2}.}
\end{equation}

Ce resultat est une regularisation UV, pas une explication de l'energie noire.
Si \(\beta_0\sim1\), la densite reste d'ordre Planck.

\section{No-go : identification directe a l'energie noire}

Si l'on impose \(\rhoreg=\rho_{\rm DE}\), alors
\begin{equation}
\boxed{
\beta_0^{(\Lambda)}
=
\left(\frac{g}{6\pi\Omega_\Lambda}\right)^{1/2}\frac{1}{H_0t_P}.
}
\end{equation}

Pour \(g=1\),
\begin{equation}
\beta_0^{(\Lambda)}=2.36\times10^{60},
\qquad
\sqrt{\beta_0^{(\Lambda)}}\lp\simeq2.5\times10^{-5}\ {\rm m}.
\end{equation}

Or cette valeur est exclue comme parametre universel par FIRAS/CMB, BBN,
fermions denses et physique atomique. Conclusion :
\begin{equation}
\boxed{
\beta_0\sim10^{60}\text{ ne peut pas etre un parametre local universel.}
}
\end{equation}

\section{Postulat III : soustraction du vide plat}

On introduit une separation entre densite brute et densite gravitationnellement
active :
\begin{equation}
\rho_{\rm brute}\neq \rhograv.
\end{equation}

Le vide plat regularise est neutralise par une prescription de renormalisation :
\begin{equation}
\boxed{T_{\mu\nu}^{\rm flat}=0.}
\end{equation}

Ce qui gravite est une difference ou une projection :
\begin{equation}
\Delta T_{\mu\nu}=T_{\mu\nu}[g]-T_{\mu\nu}[\eta],
\qquad
\boxed{
\rhograv=\mathcal P_{\rm holo}(\rhoreg).
}
\end{equation}

Ce postulat est indispensable mais non derive ici.

\section{Postulat IV : projection holographique IR}

Une region de rayon \(L\) ne peut pas contenir une energie superieure a celle
d'un trou noir de rayon \(L\) :
\begin{equation}
E(L)\leq E_{\rm BH}(L),
\qquad
E_{\rm BH}(L)=\frac{c^4L}{2G}.
\end{equation}

Avec \(V=4\pi L^3/3\), on obtient
\begin{equation}
\boxed{
\rho_{\rm BH}(L)=\frac{3c^4}{8\pi GL^2}.
}
\end{equation}

La prescription effective minimale est
\begin{equation}
\boxed{
\rhograv(L)=
\min\left[
\frac{g\rp}{16\pi^2\beta_0^2},
\frac{3c^4}{8\pi GL^2}
\right].
}
\end{equation}

Pour \(L=\LH=c/H_0\),
\begin{equation}
\rho_{\rm BH}(\LH)=\frac{3c^2H_0^2}{8\pi G}=\rhc.
\end{equation}

La densite d'energie noire observee s'ecrit donc
\begin{equation}
\rho_{\rm DE}=\Omega_\Lambda\rho_{\rm BH}(\LH).
\end{equation}

\subsection{Bits d'horizon}

On peut retrouver l'echelle
\begin{equation}
\rho_{\rm holo}\sim\frac{c^4}{GL^2}
\end{equation}
en combinant \(N_{\rm holo}\sim L^2/\lp^2\) et une energie par bit
\(\varepsilon_L\sim\hbar c/L\). Les facteurs numeriques exacts dependent de la
normalisation de \(\varepsilon_L\), de la temperature d'horizon et du comptage
des bits ; cette derivation est donc une intuition d'echelle, pas une preuve de
coefficient.

\section{Identite Planck--Hubble--energie noire}

Comme
\begin{equation}
\rp\lp^2=\frac{c^4}{G},
\end{equation}
on a
\begin{equation}
\rho_{\rm DE}
=\Omega_\Lambda\frac{3c^2H_0^2}{8\pi G}
=\boxed{
\frac{3\Omega_\Lambda}{8\pi}\rp\left(\frac{\lp}{\LH}\right)^2.
}
\end{equation}

Cette relation est une identite exacte, pas encore une dynamique.

\section{Postulat V : choix dynamique de la longueur cosmologique}

La theorie doit selectionner une longueur IR
\begin{equation}
L=L[g_{\mu\nu}].
\end{equation}

Possibilites :
\begin{equation}
L=\frac{c}{H},
\qquad
L=\RL,
\qquad
L=R_{\rm eh}=a(t)\int_t^\infty\frac{c\,\dif t'}{a(t')},
\qquad
L=R_{\rm app}.
\end{equation}

Le choix \(L=c/H(t)\) est fortement contraint dans la conjecture
\(\alpha\)--\(L\), car il induit une variation temporelle de \(\alpha\). Une
echelle quasi constante, comme \(R_\Lambda\), est donc plus naturelle.

\section{Postulat VI : tenseur energie-impulsion effectif}

Une densite correcte ne suffit pas. Il faut obtenir
\begin{equation}
T_{\mu\nu}^{\rm eff}\simeq -\rho_{\rm DE}g_{\mu\nu}.
\end{equation}

On peut ecrire formellement
\begin{equation}
S_{\rm holo}=-\int \dif^4x\sqrt{-g}\,\rho_{\rm holo}(L[g]),
\qquad
\rho_{\rm holo}(L)=\eta\frac{3c^4}{8\pi GL^2}.
\end{equation}

Le tenseur est
\begin{equation}
T_{\mu\nu}^{\rm holo}
=-\frac{2}{\sqrt{-g}}\frac{\delta S_{\rm holo}}{\delta g^{\mu\nu}}.
\end{equation}

Si \(L\) est constant,
\begin{equation}
T_{\mu\nu}^{\rm holo}=-\rho_{\rm holo}g_{\mu\nu},
\qquad
w=-1.
\end{equation}

Si \(L\) depend de la metrique, les variations de \(L[g]\) donnent des
corrections a \(w\).

\section{Prediction cosmologique : equation d'etat}

On pose
\begin{equation}
\rho_{\rm DE}(a)=\eta(a)\frac{3c^4}{8\pi G L(a)^2}.
\end{equation}

La conservation
\begin{equation}
\frac{\dif\ln\rho}{\dif\ln a}=-3(1+w)
\end{equation}
donne
\begin{equation}
\boxed{
w(a)= -1+\frac{2}{3}\frac{\dif\ln L}{\dif\ln a}
       -\frac{1}{3}\frac{\dif\ln\eta}{\dif\ln a}.
}
\end{equation}

Si \(\eta\) est constante,
\begin{equation}
\boxed{
w(a)= -1+\frac{2}{3}\frac{\dif\ln L}{\dif\ln a}.
}
\end{equation}

\section{Conjecture electromagnetique : alpha comme distance entropique}

On introduit
\begin{equation}
\calL_\Lambda=\ln\left(\frac{\RL}{10\pi\lp}\right).
\end{equation}

La conjecture est
\begin{equation}
\boxed{\alpha^{-1}=\calL_\Lambda.}
\end{equation}

Equivalent :
\begin{equation}
\boxed{
\Lambda=\frac{3}{100\pi^2\lp^2}e^{-2/\alpha}.
}
\end{equation}

\subsection{Forme entropique}

\begin{equation}
S_\Lambda=\frac{\kb\pi\RL^2}{\lp^2},
\qquad
\boxed{
\alpha^{-1}=\frac12\ln\left(\frac{S_\Lambda}{S_0^{U(1)}}\right).
}
\end{equation}

Pour reproduire le facteur \(10\pi\), il faut
\begin{equation}
S_0^{U(1)}=100\pi^3\kb.
\end{equation}

Une parametrisation est
\begin{equation}
S_0^{U(1)}=4\pi\left[2\pi g_{\rm eff}^{U(1)}\right]^2\kb,
\qquad
g_{\rm eff}^{U(1)}=\frac52=2+\frac12.
\end{equation}

Le verrou est de prouver
\begin{equation}
\boxed{
g_{\rm edge}^{U(1)}=\frac12
}
\end{equation}
comme capacite holographique renormalisee du mode de bord Maxwell.

\section{Relations induites par la conjecture alpha}

Si \(\alpha=1/\calL_\Lambda\), alors
\begin{align}
e&=\frac{q_P}{\sqrt{\calL_\Lambda}},
\qquad
q_P=\sqrt{4\pi\varepsilon_0\hbar c},\\
v_e&=\alpha c=\frac{c}{\calL_\Lambda},\\
a_0&=\frac{\hbar}{m_ec\alpha}=\bar\lambda_e\calL_\Lambda,\\
r_e&=\alpha\bar\lambda_e=\frac{\bar\lambda_e}{\calL_\Lambda},\\
E_R&=\frac12m_ec^2\alpha^2=\frac{m_ec^2}{2\calL_\Lambda^2},\\
\frac{F_e}{F_G}&=\frac{m_P^2}{m_e^2\calL_\Lambda}.
\end{align}

Ces relations sont des reecritures algebriques de formules connues. Leur
interet depend entierement de la validite de la conjecture \(\alpha\)--\(R_\Lambda\).

\section{Position--vitesse et electron atomique}

En regime non relativiste,
\begin{equation}
\dif N=\frac{g m^3\dif^3x\,\dif^3v}{h^3(1+\beta m^2v^2)^3}.
\end{equation}

La cellule effective est
\begin{equation}
\boxed{
\Delta^3x\,\Delta^3v\sim
\left(\frac{h}{m}\right)^3(1+\beta m^2v^2)^3.
}
\end{equation}

Le seuil cinematique est
\begin{equation}
v_{\rm GUP}=c\frac{m_P}{m\sqrt{\beta_0}},
\qquad
\Delta x_{\min}=\sqrt{\beta_0}\lp.
\end{equation}

Pour l'electron atomique,
\begin{equation}
\chi_e=\beta m_e^2v_e^2
=\beta_0\left(\frac{m_e}{m_P}\right)^2\frac{1}{\calL_\Lambda^2}.
\end{equation}

Pour \(\beta_0\sim1\), cette correction est minuscule. Pour
\(\beta_0\sim10^{60}\), elle serait catastrophique, ce qui renforce le no-go
d'un \(\beta_0\) universel tres grand.

\section{Extension aux autres interactions}

Une generalisation speculative serait
\begin{equation}
\boxed{
\alpha_i^{-1}=\ln\left(\frac{\RL}{4\pi g_{\rm eff}^{(i)}\lp}\right).
}
\end{equation}

Mais les constantes de jauge courent avec l'echelle de renormalisation. Pour
le modele standard, il faudrait traiter \(U(1)_Y\), \(SU(2)_L\), \(SU(3)_c\),
le melange electroweak et le confinement QCD. Cette extension est donc un
programme, pas un resultat.

\section{Six blocs ouverts}

\begin{longtable}{@{}p{3.5cm}p{11cm}@{}}
\toprule
Bloc & Objectif \\
\midrule
1. Covariance GUP & Construire une mesure covariante compatible avec FLRW. \\
2. Vide plat & Deriver la neutralisation de \(T_{\mu\nu}^{\rm flat}\). \\
3. Projection IR & Deriver dynamiquement la projection holographique. \\
4. Longueur \(L\) & Choisir \(L[g_{\mu\nu}]\) par un principe causal. \\
5. Tenseur effectif & Obtenir \(T_{\mu\nu}\simeq-\rho g_{\mu\nu}\). \\
6. Prediction \(w(a)\) & Produire une fonction \(w(a)\) testable. \\
\bottomrule
\end{longtable}

\section{Falsifiabilite}

Le cadre serait affaibli si :
\begin{itemize}
\item la relation \(\alpha^{-1}=\ln(\RL/10\pi\lp)\) echoue avec des parametres cosmologiques affines ;
\item le facteur \(10\pi\) ne peut pas etre relie a un seuil \(U(1)\) naturel ;
\item la projection holographique ne peut pas produire un tenseur d'energie sombre ;
\item la mesure GUP ne peut pas etre covariantisee.
\end{itemize}

\section{Architecture finale}

\subsection{Chaine cosmologique}

\begin{equation}
\boxed{
\text{cellule GUP}\Rightarrow \rhoreg<\infty
\Rightarrow \text{projection holographique}
\Rightarrow \rho_{\rm DE}.
}
\end{equation}

\begin{equation}
\frac{g\dif^3x\dif^3p}{h^3(1+\beta p^2)^3}
\Rightarrow
\frac{g\rp}{16\pi^2\beta_0^2}
\Rightarrow
\min\left[\rhoreg,\frac{3c^4}{8\pi GL^2}\right].
\end{equation}

\subsection{Chaine electromagnetique}

\begin{equation}
\boxed{
\text{horizon de de Sitter}+\text{bord }U(1)
\Rightarrow \alpha.
}
\end{equation}

\begin{equation}
S_\Lambda=\frac{\kb\pi\RL^2}{\lp^2},
\qquad
S_0^{U(1)}=4\pi\left[2\pi\left(2+\frac12\right)\right]^2\kb,
\qquad
\alpha^{-1}=\frac12\ln\left(\frac{S_\Lambda}{S_0^{U(1)}}\right).
\end{equation}

\section{Conclusion}

La theorie holographique UV/IR a cellule de phase deformee repose sur une idee :
\begin{equation}
\boxed{
\begin{gathered}
\text{les constantes observables ne sont pas fixees par l'UV seul,}\\
\text{mais par une tension entre l'UV quantique et l'IR cosmologique.}
\end{gathered}
}
\end{equation}

Le GUP fournit une densite locale finie d'etats. La gravite impose une
projection holographique globale. Le secteur \(U(1)\) de bord pourrait fixer
la constante de structure fine par une distance entropique entre l'horizon de
de Sitter et un seuil electromagnetique minimal.

Le cadre n'est pas confirme. Son interet est de reduire plusieurs problemes a
des verrous precis :
\begin{equation}
\boxed{
\text{covariance GUP, projection IR dynamique, choix de }L, T_{\mu\nu}, w(a),
\text{ et bord }U(1).
}
\end{equation}

\begin{thebibliography}{99}

\bibitem{CODATA2022}
P. J. Mohr, D. B. Newell, B. N. Taylor and E. Tiesinga,
``CODATA recommended values of the fundamental physical constants: 2022'',
\emph{Reviews of Modern Physics} \textbf{97}, 025002 (2025).

\bibitem{Planck2018}
N. Aghanim et al. (Planck Collaboration),
``Planck 2018 results. VI. Cosmological parameters'',
\emph{Astronomy \& Astrophysics} \textbf{641}, A6 (2020), arXiv:1807.06209.

\bibitem{CKN}
A. Cohen, D. Kaplan and A. Nelson,
``Effective Field Theory, Black Holes, and the Cosmological Constant'',
\emph{Physical Review Letters} \textbf{82}, 4971 (1999).

\bibitem{LiHDE}
M. Li,
``A Model of Holographic Dark Energy'',
\emph{Physics Letters B} \textbf{603}, 1--5 (2004).

\bibitem{Kempf}
A. Kempf, G. Mangano and R. B. Mann,
``Hilbert space representation of the minimal length uncertainty relation'',
\emph{Physical Review D} \textbf{52}, 1108 (1995).

\bibitem{DonnellyWall}
W. Donnelly and A. C. Wall,
``Entanglement entropy of electromagnetic edge modes'',
\emph{Physical Review Letters} \textbf{114}, 111603 (2015).

\bibitem{BallLawWong}
A. Ball, Y. T. A. Law and G. Wong,
``Dynamical Edge Modes and Entanglement in Maxwell Theory'',
arXiv:2403.14542.

\end{thebibliography}

\end{document}
